\section{Problem Definition and Threat Model}
\label{sec:model}

We model an ordinary embodied VLA policy call and define the evidence required
before its output may affect task state.  The formulation is independent of the
evaluation platform; Section~\ref{sec:evaluation} later instantiates it in
LIBERO-Safety.

\subsection{System Model and Protected Object}

Let $x_t$ denote the embodied system state at epoch $t$, $T$ an authoritative
task artifact, and $O_t^T$ the observation used by the monitor.  The VLA
receives a policy view $(P_t^{pol},O_t^{pol},H_t^{pol})$ containing its prompt,
observation, and history.  These views may agree during nominal execution or
diverge after input corruption.  Their separation is a TCB assumption: logging
policy-facing bytes does not make those bytes authoritative.

At state epoch $t$, the two branches compute
\begin{equation}
\begin{aligned}
Z_t &= \mathsf{SelectFrozen}(T,O_t^T),\\
A_t &= \pi(P_t^{pol},O_t^{pol},H_t^{pol}),\\
S_t &= \mathsf{AssessLocal}(Z_t,O_t^T,A_t).
\end{aligned}
\end{equation}
$Z_t$ is selected and frozen before the policy returns $A_t$.  It is an
independent monitor anchor, not a recovered latent intent or a hidden prompt to
the policy.  The compiler maps $T$ to a trusted task graph
$\mathcal G_T=(V_T,E_T,\ell_T,\Phi_T)$: vertices are legal task phases, edges
are permitted phase transitions, $\ell_T$ gives each vertex a canonical
skill/entity tuple, and $\Phi_T(O_t^T)\subseteq V_T$ is the state-enabled
frontier.  \textsf{SelectFrozen} may emit $Z_t=\ell_T(v_t)$ only for
$v_t\in\Phi_T(O_t^T)$; otherwise it emits \texttt{unknown}.  The graph encodes
semantic authority, not a reference trajectory.  In deployment, $T$ could be a
signed job ticket or protected workflow state.  Without an independent task
source, \system cannot create authority and its semantic claim does not apply.

The protected object is the exact online block
$A_t\in\mathbb{R}^{H\times d}$, represented by its ordered numeric values and
shape in the runtime's canonical domain.  Assessment $S_t$ binds $A_t$ to
$(Z_t,O_t^T)$.  The consumer then creates execution contract $C_t$, one-use
authorization $\mathit{Auth}_t$, ordered receipts $R_t$, and effect evidence
$E_t$; $g_t$ denotes the selected execution configuration.  Every object binds
the episode nonce, proposal index, and state epoch; changes invalidate reuse.
The command identity denotes equality in a canonical numeric domain, not raw
memory equality or a guaranteed physical trajectory.  The executable identity
is at least $(A_t,g_t,t)$ because the same numeric block can realize differently
under a different controller or guard configuration.

\subsection{Attack Path and Adversary}

The adversary controls inputs and software outside the TCB.  Its path begins
after the secure split: it may modify the policy-facing instruction,
digital/software visual input, or policy history while leaving the protected
task branch unchanged.  The resulting frozen-policy call can produce a
task-divergent or risk-inducing source block.  After assessment, the adversary
may attempt command substitution, step reordering, stale-authorization replay,
cross-proposal receipt/effect splicing, or omission of unfavorable evidence.
Its objective is either to cause an unauthorized or covered-risk dispatch, or
to make a failed transaction appear eligible for task-phase advance.

This model does not assume that the policy checkpoint is malicious: a known,
immutable policy remains an untrusted proposer because its context is
attackable.  Conversely, controlling code outside the TCB does not grant the
adversary a second actuator path; complete mediation by the registered
dispatch boundary is a trusted assumption stated below.  Our empirical attack
study instantiates instruction modification, and the focused integrity suite
instantiates the modeled transaction faults.  Visual and history attacks are
in system scope but are not empirically evaluated.

\subsection{Trusted Assumptions and Evidence Endpoints}

The TCB contains the authoritative task source, trusted observation tap,
task-graph compiler and selector, local checker, contract compiler,
nonce/digest ledger, single dispatch boundary, effect observer, and protected
evidence storage.  The policy, policy-facing prompt/observation/history, and
source \actionblock are untrusted.  This is a software reference-monitor model:
we assume the TCB code and its cryptographic primitives execute as specified;
the trusted tap precedes the modeled policy-view transformation; every
actuator command traverses the single dispatch boundary; and receipt/effect
storage cannot be rewritten outside the TCB.

Three mediated endpoints delimit the evidence chain: the authority tap captures
$T$ and $O_t^T$ before policy-view modification; the dispatch boundary records
the command accepted by the actuator interface; and the effect window opens
only after dispatch.  Evidence outside these endpoints cannot establish task
authority, command application, or a post-dispatch effect, respectively.  In
particular, a policy explanation shares the attacked context, a pre-dispatch
prediction is not a receipt, and an unbound later observation is not evidence
for the current proposal.

\subsection{Proof Alignment Property}

Collect the evidence for one proposal as
$\chi_t=(T,O_t^T,Z_t,A_t,S_t,C_t,g_t,\mathit{Auth}_t,R_t,E_t)$.  We define the
top-level checker-relative property
\begin{equation}
\begin{aligned}
\mathsf{ProofAligned}(\chi_t)\triangleq {}&
\mathsf{Eligible}_t\land\mathsf{TxnAligned}_t\\
&\land(\mathsf{Triggered}_t\Rightarrow\mathsf{Contained}_t).
\end{aligned}
\label{eq:proof-aligned}
\end{equation}
The three conjuncts correspond directly to the three challenges in
Section~\ref{sec:background}.  They are defined below in terms of evidence that
the runtime can produce and check.

\textbf{P1---checker-relative eligibility.}  A block may receive covered-risk
authorization only if a qualified checker assessed that exact block against a
legal $Z_t$ and the matching trusted observation epoch.  The implemented
condition is

\begin{align}
\mathsf{Eligible}_t \triangleq {}&
  \mathsf{TrustedProv}(T,O_t^T,Z_t) \nonumber\\
&\land \mathsf{LegalFrontier}(T,O_t^T,Z_t) \nonumber\\
&\land \mathsf{Bound}(S_t,Z_t,O_t^T,A_t) \nonumber\\
&\land \mathsf{ScreenAvailable}(S_t) \nonumber\\
&\land \neg\mathsf{HardRisk}(S_t) \nonumber\\
&\land \mathsf{Qualified}(S_t.\mathit{assessor}).
\end{align}

This definition establishes provenance, frontier legality, assessment identity,
and covered hard gates.  It does not establish that every admissible action
advances $T$.  Task-progress and unavailable articulation evidence are
advisory and force next-block replanning; velocity, workspace, unexpected
contact, stale/malformed commands, and unsupported unknown evidence are hard
gates.

\textbf{P2---transaction alignment.}  A task phase may advance only if the
authorization is fresh and bound to the proposal, the dispatched prefix is
ordered and exact, and the post-dispatch evidence satisfies the contract:
\begin{align}
\mathsf{TxnAligned}_t \triangleq {}&
  \mathsf{FreshAndBound}(C_t,\mathit{Auth}_t,R_t) \nonumber\\
&\land \mathsf{OrderedExactPrefix}(R_t,\mathit{Auth}_t) \nonumber\\
&\land \mathsf{AfterDispatch}(E_t,R_t) \nonumber\\
&\land \mathsf{Required}(C_t)\subseteq E_t \nonumber\\
&\land \mathsf{Forbidden}(C_t)\cap E_t=\emptyset \nonumber\\
&\land \neg\mathsf{ObsViolation}(E_t).
\end{align}
The contract binds the action, semantic-subtask, policy-prompt, assessment, and
trusted-observation digests; state epoch; required/forbidden effects; and the
observation window.  An authorization can be consumed once.  In Python,
$\mathsf{OrderedExactPrefix}$ requires contiguous receipt indices and checks
each applied-step digest against \texttt{authorization.action\_at(i)}.  The
Lean model represents the same ordered per-step digest list and proves that
every bound receipt resolves at its index to the applied digest; Python
additionally enforces contiguous indices.  A complete serializer refinement is
outside the theorem scope.  An incomplete prefix, unknown effect, or receipt
from another proposal cannot advance the phase.

\textbf{P3---covered physical containment.}  Let $g_t$ be the selected guard
configuration and $m_j(g_t,A_t,O_t^T)$ the predicted next-step margin for joint
side $j$.  A triggered screen qualifies only when
\begin{align}
\mathsf{Contained}_t \triangleq {}&
  \mathsf{SameSource}(g_t,A_t)\land\mathsf{Restored}(g_t) \nonumber\\
&\land \min_{j\in\{1,\ldots,14\}}m_j(g_t,A_t,O_t^T)\geq0.15 \nonumber\\
&\land \mathsf{Force}(g_t,A_t,O_t^T)\leq 10{,}000.
\end{align}
The goal is covered joint-side containment, not exact trajectory reproduction
or arbitrary collision avoidance.

The distinction between source action and executable configuration is
important.  A virtual guard leaves $A_t$ unchanged but changes how the
controller maps its values to motion.  The runtime records this configuration,
but the current Lean contract does not yet type-bind it as an independent
digest.  We expose that refinement gap rather than equating command identity
with physical-trajectory identity.

Task-state progress is stricter than local acceptance.  Let $q_t$ and
$q_{t+1}$ be consecutive task phases.  The monitor enforces
\begin{equation}
\begin{aligned}
\mathsf{Advance}(q_t,q_{t+1})
&\Rightarrow \mathsf{ProofAligned}(\chi_t)\\
&\phantom{\Rightarrow{}}\land
\mathsf{CompletionObserved}(q_t,q_{t+1},E_t).
\end{aligned}
\label{eq:phase-advance}
\end{equation}
Lean mechanizes the finite identity, authorization, receipt, evidence, and
phase-transition subrelation of Eqs.~\ref{eq:proof-aligned}
and~\ref{eq:phase-advance}; task selection, checker soundness, and continuous
containment remain explicit qualification obligations.

\subsection{Out of Scope}

We exclude task-source or monitor compromise, attacks that corrupt the monitor
observation before the trusted tap, forged actuator feedback inside the TCB,
and actuator paths that bypass the mediator.  Hashes, nonces, and Lean semantics
do not supply a hardware root of trust, sensor correctness, arbitrary-dynamics
robustness, or hard-real-time guarantees.  If the trusted and policy views share
a compromised source, or the observer omits a relevant effect, a transaction
may be internally aligned yet wrong about the physical world.  The empirical
scope of the current implementation is stated separately in
Section~\ref{sec:evaluation} and Section~\ref{sec:discussion}.
